% !PRINTED 515
\weberpage{515}

Das Product (4) gestattet nun die Substitutionen der Gruppe $Q$,
und wenn also $\psi$ wie oben eine zu dieser Gruppe gehörige
Function in $N$ bedeutet, so lässt sich nach \S.\,155, 2. die
Function $g_1(t,\varepsilon)$ rational durch $\psi$ ausdrücken,
d.\,h. $g_1(t,\varepsilon)$ ist eine Function von $t$ im Körper
$j^{\mathrm{ten}}$ Grades $\Omega(\psi)$, und soll demnach durch
$g(t,\psi)$ bezeichnet werden.

Die Grösse $\psi$ ist eine der Wurzeln einer irreducibeln
Gleichung vom Grade $j$

\weq{6}{\varphi(u) = 0,}

deren übrige Wurzeln $\psi_1, \psi_2 \dots \psi_{j-1}$ sind, und
nach \S.\,155, (6) ist $g(t)$ in die Factoren zerlegbar

\weq{7}{g(t) = g(t,\psi)\, g(t,\psi_1) \dots g(t,\psi_{j-1}).}

Nun ist

\[ g(t,\psi) = g_1(t,\varepsilon), \]

und aus (7) folgt, dass diese Gleichung nicht bestehen bleibt,
wenn auf der linken Seite für $\psi$ eine der anderen Wurzeln
von (6) gesetzt wird. Man kann also für $t$ einen solchen
rationalen Werth setzen, dass die Gleichung

\weq{8}{g(t,u) - g_1(t,\varepsilon) = 0}

nur die eine Wurzel $u = \psi$ mit der Gleichung (6) gemein hat.
Der grösste gemeinschaftliche Theiler von (6) und (8) ist dann
in Bezug auf $u$ linear, und wenn man darin $u = \psi$ setzt, so
ergiebt sich, \gesperrt{dass $\psi$ rational durch $\varepsilon$
ausdrückbar ist}. Der Körper $\Omega(\psi)$ ist also sicher ein
Theiler des Körpers $\Omega(\varepsilon)$. Der Grad des Körpers
$\Omega(\varepsilon)$, d.\,h. der Grad der Hülfsgleichung
$\chi(\varepsilon) = 0$, ist daher ein Vielfaches von $j$ und
niemals niedriger als $j$. Wenn der Grad der Hülfsgleichung
gleich $j$ ist, was z.\,B. dann immer eintritt, wenn der Grad
von $\chi$ eine Primzahl ist, so ist $\Omega(\varepsilon)$ mit
$\Omega(\psi)$ identisch, d.\,h. $\varepsilon$ ist auch rational
durch $\psi$ ausdrückbar. In diesem Falle sind die conjugirten
Werthe $\varepsilon, \varepsilon_1 \dots \varepsilon_{j-1}$ die
sämmtlichen Wurzeln von $\chi = 0$, und man erhält sie, wenn man
in dem rationalen Ausdruck von $\varepsilon$ durch $\psi$ die
Function $\psi$ durch jede der conjugirten Functionen $\psi,
\psi_1 \dots \psi_{j-1}$ ersetzt. Es kann also $\varepsilon$
selbst für $\psi$ genommen werden, und man findet eine
Zerlegung

\[ g(t) = g(t,\varepsilon)\, g(t,\varepsilon_1) \dots g(t,\varepsilon_{j-1}). \]

Die Grössen $\varepsilon, \varepsilon_1, \varepsilon_2 \dots
\varepsilon_{j-1}$ gehören sämmtlich dem Körper $\Omega(\varrho)$
an.
